\section{Certified stopping with backtracking}
\label{sec:stopping}
Assume the convex model of Section~\ref{sec:model}. Convexity gives
\begin{equation}\label{eq:gap_certificate}
 f(x)-f(x^*)\le
 \langle\nabla f(x),x-x^*\rangle\le G(x).
\end{equation}
Thus the full gap is a computable error certificate even when
$\nabla f(x^*)\ne0$. No estimate of that unknown gradient is needed.
Use a positive global curvature bound $K_Q$ from
\eqref{eq:global_curvature} in the analysis; if that bound is zero,
it may be enlarged to any positive value. The algorithm does not
need $K_Q$, $L_0$, or $L_1$ as inputs.

\paragraph{Exact-oracle algorithm.}
Fix $\varepsilon>0$ and an initial stored curvature $L_{\rm init}>0$.
At a visited point $x$, compute its gradient, a full linear minimizer
$s\in Q$, and $G(x)$. Retain
\[
 \ell\leftarrow\max\{\ell,f(x)-G(x)\},
\]
where the first query initializes $\ell=f(x_0)-G(x_0)$.
Stop if $f(x)-\ell\le\varepsilon$. Otherwise set $d=s-x$, start with
half the previous accepted curvature (initially $L_{\rm init}$),
and double $L$ until
\begin{equation}\label{eq:stopping_bt}
 \alpha=\min\left\{1,\frac{G(x)}{L\|d\|^2}\right\},\qquad
 f(x+\alpha d)\le f(x)-\alpha G(x)+\frac{L\alpha^2\|d\|^2}{2}.
\end{equation}
Accept $x+\alpha d$, cache its objective value, and store $L$.
The update is ordinary Frank-Wolfe backtracking \cite{pedregosa}.
Objective values, gradients, and full linear minimizers are exact.

\begin{theorem}[Finite stopping and complete oracle counts]
\label{thm:finite_stopping}
Every line search terminates. The algorithm stops with
$f(x_N)-f(x^*)\le f(x_N)-\ell_N\le\varepsilon$.
If $G(x_0)>\varepsilon$, set
\[
 J=\left\lceil\log_2\frac{G(x_0)}{\varepsilon}\right\rceil,
 \quad b_j=G(x_0)/2^j\ (0\le j<J),
\]
\[
 H_{\rm init}=\left\lceil
 \log_2\max\left\{1,\frac{L_{\rm init}}{2K_Q}\right\}\right\rceil.
\]
The number $N$ of accepted updates satisfies
\begin{align}
 N&\le H_{\rm init}+\left\lceil\sum_{j=0}^{J-1}
       \max\left\{4,\frac{16K_QD^2}{b_j}\right\}\right\rceil
       \label{eq:stopping_steps}\\
  &\le H_{\rm init}+\left\lceil4J+
       \frac{32K_QD^2}{\varepsilon}\right\rceil.
       \label{eq:stopping_simple}
\end{align}
There are exactly $N+1$ gradient evaluations and $N+1$ full linear
minimization calls, including the terminal query. The initial objective
value and all accepted and rejected line-search trials number at most
\begin{equation}\label{eq:stopping_values}
 N_f\le 1+2N+
 \left\lceil\log_2\max\left\{1,\frac{2K_Q}{L_{\rm init}}\right\}\right\rceil.
\end{equation}
If $G(x_0)\le\varepsilon$, there are no updates and exactly one call
of each oracle; the phase quantities are unnecessary.
\end{theorem}

\begin{proof}
Every retained $\ell$ is a lower bound on $f(x^*)$, so
\[
 0\le f(x)-f(x^*)\le f(x)-\ell\le G(x).
\]
A nonterminal point therefore has $G(x)>0$ and $d\ne0$.
Every $L\ge K_Q$ passes \eqref{eq:stopping_bt}. Accepted values obey
\begin{equation}\label{eq:stopping_progress}
 \begin{split}
 L_k&\le\max\{L_{k-1}/2,2K_Q\},\qquad L_{-1}=L_{\rm init},\\
 f(x_k)-f(x_{k+1})&\ge
 \min\left\{\frac{G(x_k)}2,\frac{G(x_k)^2}{2L_kD^2}\right\}.
 \end{split}
\end{equation}
Thus the objective decreases, $\ell$ increases, and $f(x_k)-\ell_k$
decreases. If $L_k>2K_Q$, the accepted trial was the first: otherwise
its preceding rejected value $L_k/2>K_Q$ would have passed.
Such an update halves the inherited curvature. Once $L_k\le2K_Q$,
it stays at most $2K_Q$. There are at most $H_{\rm init}$ updates
above this threshold.

Assign every other nonterminal update to
$b_j/2<f(x_k)-\ell_k\le b_j$. These phases are visited in order.
Here $G(x_k)>b_j/2$, and \eqref{eq:stopping_progress} gives decrease
at least
\[
 \min\left\{b_j/4,\ b_j^2/(16K_QD^2)\right\}.
\]
The total decrease of all updates starting in that phase is at most
the true error at its entry, hence at most $b_j$. This includes the
last update leaving the phase, because its endpoint has objective
at least $f(x^*)$. Division and summation prove
\eqref{eq:stopping_steps}. The inequality
$\sum_{j<J}b_j^{-1}<2/\varepsilon$ gives
\eqref{eq:stopping_simple}. The same finite bound applies to every
unfinished prefix, excluding an infinite nonterminal run.

If update $k$ rejects $q_k$ trials, then
$L_k=2^{q_k-1}L_{k-1}$. For $N\ge1$,
\[
 \sum_{k<N}(q_k+1)=2N+\log_2\frac{L_{N-1}}{L_{\rm init}},
 \quad L_{N-1}\le\max\{L_{\rm init}/2^N,2K_Q\}.
\]
Adding the initial objective value proves \eqref{eq:stopping_values}.
Each visited point uses one gradient and one full linear minimizer.
\end{proof}

\subsection{Carrying absolute directional accuracy between iterations}
A directional query with requested accuracy $\eta>0$ returns $h\ge0$
and $d=s-x$, $s\in Q$, with certified bounds
\begin{equation}\label{eq:directional_query}
 G(x)\le h+\eta,\qquad
 -\langle\nabla f(x),d\rangle\ge h-\eta.
\end{equation}
This is an explicit oracle requirement, not a consequence of the line-search
check. Objective values remain exact. The exact algorithm above needs
no such additional oracle.

Start with $\eta_0>0$ and fix $0<\theta<1$. After every query, update
$\ell\leftarrow\max\{\ell,f(x)-(h+\eta)\}$ and first test
$f(x)-\ell\le\varepsilon$. If this fails and $\eta>\theta h$,
halve $\eta$ and query again at the same point, reusing its objective
value. Once $\eta\le\theta h$, use \eqref{eq:stopping_bt} with
$G(x)$ replaced by $h-\eta$. Carry the current $\eta$ to the next point.

\begin{corollary}[All accuracy refinements are counted]
\label{cor:directional_stopping}
The rule stops with $f(x_N)-f(x^*)\le\varepsilon$.
Put $\kappa=(1-\theta)/(1+\theta)$. The total number $R$ of accuracy
halvings obeys
\begin{equation}\label{eq:accuracy_refinements}
 R\le\left\lceil\log_2\max\left\{1,
 \frac{(1+\theta)\eta_0}{\theta\varepsilon}\right\}\right\rceil.
\end{equation}
If the first upper gap $h_0+\eta_0>\varepsilon$, define $J$ and $b_j$
as above with $G(x_0)$ replaced by $h_0+\eta_0$. Then
\begin{equation}\label{eq:directional_steps}
 N\le H_{\rm init}+\left\lceil
 \frac{4J}{\kappa}+\frac{32K_QD^2}{\kappa^2\varepsilon}
 \right\rceil.
\end{equation}
There are exactly $N+1+R$ directional queries. The objective bound
\eqref{eq:stopping_values} is unchanged. If each directional query uses
one full linear minimization call, their number is also $N+1+R$.
If the first upper gap is at most $\varepsilon$, the rule stops immediately.
\end{corollary}
\begin{proof}
At a failed nonterminal accuracy check,
\[
 \varepsilon<f(x)-\ell\le h+\eta
 <(1+\theta)\eta/\theta.
\]
Since $\eta$ is halved and never reset, this proves
\eqref{eq:accuracy_refinements}. At a passed nonterminal check,
\[
 h-\eta\ge\kappa(h+\eta)\ge\kappa(f(x)-\ell)>0.
\]
The preceding descent and phase proof therefore applies with phase
slope threshold $\kappa b_j/2$, giving
\eqref{eq:directional_steps}. Each visited point has one closing query,
including the terminal point, and each halving adds one failed query.
Refinements at a fixed point do not evaluate the objective again.
\end{proof}

The stopping proof uses the full gap and permits a nonzero limiting
gradient. In Section~\ref{sec:normal}, the true curvature bound and the
complete execution are independent of an arbitrarily large normal
gradient. Directional-query counts specify how often the requested
accuracy is needed; the cost of producing each answer is determined by
the chosen oracle implementation.
